\documentclass{article}
\usepackage{amsmath, amssymb, amsthm}
\usepackage{hyperref}
\usepackage{geometry}
\geometry{margin=1in}

\title{Erd\H{o}s Problem \#601}
\date{Last updated: 2025-08-31}
\author{Source: erdosproblems.com}

\begin{document}
\maketitle

\noindent\textbf{Status:} OPEN \quad \textbf{Prize: \$500}

\noindent\textbf{Tags:} graph theory, set theory

\medskip
\noindent\textbf{Problem Statement:}

For which limit ordinals $\alpha$ is it true that if $G$ is a graph with vertex set $\alpha$ then $G$ must have either an infinite path or independent set on a set of vertices with order type $\alpha$?

\bigskip
\noindent\textbf{Remarks:}

A problem of Erd\H{o}s, Hajnal, and Milner \cite{EHM70}, who proved this is true for $\alpha < \omega_1^{\omega+2}$. 

In \cite{Er82e} Erd\H{o}s offers \$250 for showing what happens when $\alpha=\omega_1^{\omega+2}$ and \$500 for settling the general case.

Larson \cite{La90} proved this is true for all $\alpha<2^{\aleph_0}$ assuming Martin's axiom.

\bigskip
\noindent\textbf{References:}

[EHM70] Erd\H{o}s, P. and Hajnal, A. and Milner, E. C., Set mappings and polarized partition relations. Combinatorial theory and its applications, I-III (Proc.
Colloq., Balatonf\"{u}red, 1969) (1970), 327-363.
 
 [Er82e] Erd\H{o}s, Paul, Some of my favourite problems which recently have been solved. (1982), 59--79.
 
 [La90] Larson, Jean A., Martin's axiom and ordinal graphs: large independent sets or infinite paths. Ann. Pure Appl. Logic (1990), 31-39.

\bigskip
\noindent\small{Source: \url{https://www.erdosproblems.com/601}}
\end{document}
